% Generated draft: requires mathematical and source review.
\documentclass[11pt]{article}
\usepackage[margin=25mm]{geometry}
\usepackage{fontspec}
\setmainfont{Latin Modern Roman}
\usepackage{amsmath,amssymb,mathtools,unicode-math}
\setmathfont{Latin Modern Math}
\usepackage{enumitem,needspace,xurl,hyperref}
\setlength{\parindent}{0pt}
\setlength{\parskip}{5pt}
\setlength{\emergencystretch}{3em}
\newcommand{\N}{\mathbb{N}}
\newcommand{\Z}{\mathbb{Z}}
\newcommand{\Q}{\mathbb{Q}}
\newcommand{\R}{\mathbb{R}}
\newcommand{\C}{\mathbb{C}}
\newcommand{\F}{\mathbb{F}}
\newcommand{\A}{\mathbb{A}}
\newcommand{\PP}{\mathbb{P}}
\newcommand{\E}{\mathbb{E}}
\newcommand{\eps}{\varepsilon}
\newcommand{\dd}{\,\mathrm{d}}
\begin{document}
\section*{UnsolvedMath ID: 30001442}
\textbf{Source problem code:} catalogue entry
\subsection{Definitions and mathematical statement}
Let \(X\) be a topological space. It is second-countable if its topology has a countable base, and it is \(T_0\) if any two distinct points are distinguished by an open set. A nonempty closed subset is irreducible if it is not the union of two proper closed subsets. The space is sober if every nonempty irreducible closed subset is the closure of exactly one point.

For an explicit interpretation of the non-Hausdorff terminology, call a subset compact if every open cover has a finite subcover, without imposing a Hausdorff condition, and saturated if it is an intersection of open subsets. Interpret coherent as requiring the intersection of any two compact saturated subsets to be compact. Interpret locally compact as requiring that for every \(x\in X\) and every open neighborhood \(U\) of \(x\), there exist an open set \(V\) and a compact subset \(K\) with
\[
x\in V\subseteq K\subseteq U.
\]
The supplied record does not define coherence or non-Hausdorff local compactness; these conventions require confirmation against the cited report.

For a complex \(C^*\)-algebra \(A\), let \(\operatorname{Prim}(A)\) be the set of kernels of nonzero irreducible \(*\)-representations of \(A\), equipped with the hull--kernel topology. Its closed sets are
\[
\operatorname{hull}(I)=\{P\in\operatorname{Prim}(A):I\subseteq P\},
\]
where \(I\) ranges over the closed two-sided ideals of \(A\). Interpret amenability in the Banach-algebra sense: for every Banach \(A\)-bimodule \(M\), every bounded derivation \(D:A\to M^*\) is inner. Thus a bounded linear map satisfying \(D(ab)=a\cdot D(b)+D(a)\cdot b\) must have the form \(D(a)=a\cdot\varphi-\varphi\cdot a\) for some \(\varphi\in M^*\), using the induced dual module actions.

\textbf{Question (attributed to Kirchberg by the supplied curation).} Is it true that for every second-countable, coherent, locally compact, sober \(T_0\) space \(X\), there exist an amenable \(C^*\)-algebra \(A\) and a homeomorphism
\[
h:X\longrightarrow\operatorname{Prim}(A)?
\]
No separability or unitality condition on \(A\), and no Hausdorff or global compactness condition on \(X\), is imposed by the supplied statement. The empty space is not explicitly excluded. The question remains subject to confirmation of the terminology specified above.
\subsection{Short English statement}
Does every second-countable, coherent, locally compact, sober \(T_0\) space occur as the primitive ideal space of an amenable \(C^*\)-algebra? The supplied curation attributes the question to Kirchberg and labels it open. The non-Hausdorff compactness terminology needs confirmation against the report.
\subsection{Sources}
\begin{itemize}
\item Local UnsolvedMath registry: permanent ID \(30001442\), title \emph{Realization of Sober Spaces as Primitive Ideal Spaces}, source problem code \texttt{catalogue entry}, in the Oberwolfach Reports Open Problems collection. The supplied proposed year is 2010; the proposer field is empty, while the supplied literature assessment attributes the question to Kirchberg.
\item \emph{C*-Algebren} (2010), Oberwolfach Reports, pp.\ 667--732, DOI \texttt{10.4171/owr/2010/13}. Supplied source URL and literature-source link, labeled ``original OWR report'': \url{https://doi.org/10.4171/owr/2010/13}. This report citation is supplied provenance and has not been independently inspected here.
\item Supplied literature assessment dated \texttt{2026-08-21}: the record labels the problem open and states that the report gives reductions and special cases, but that no general resolution was found in the supplied review. This is supplied dated curation, not an independently verified literature assessment or a claim about current status.
\item The supplied statement status is \texttt{corrected\_verified}; its accompanying verification note states only that TeX notation was enclosed in math delimiters without changing the mathematical content. That note does not establish mathematical correctness or verification against the report.
\end{itemize}
\end{document}
